\section{Layered scoring and scientific completion}
\label{sec:si-scoring}

\subsection{Composite and reporting layers}
For a formally organized run, the reported composite combines graph structure $G$, physical self-consistency $C$, and accuracy $A$:
\begin{equation}
S=0.1G+0.3C+0.6A.
\label{eq:si-composite}
\end{equation}
The three components are reported separately, alongside execution status, original graph delivery, and completion of the requested steps. Complete scientific validation additionally requires every requested species and step to meet its physical criteria.

\subsection{Graph structure and chemical correspondence}
The graph component measures how well the delivered mechanism covers the reference species and steps. For comparison, transition-state nodes are folded into their elementary-step edges. Species are matched by chemical identity; steps are compared by their endpoints and atom-mapped bond changes. RDKit represents the species and NetworkX represents and matches the graphs.\citep{ref13,networkx2008} Missing reference species or steps, incorrect chemical identities, and mismatched accessibility declarations add edit costs. The graph component permits additional nodes and edges.

Graph comparison ignores stereochemistry for Easy and Normal tasks and retains it for Hard tasks. A species' linked native calculation results determine which evidence is assessed for it. Graph coverage receives its own score independently of the binary self-consistency result.

\subsection{Graded physical self-consistency}
Self-consistency is scored on the species and steps in the organized graph, independently of numerical agreement with Gold. Table~\ref{tab:si-self-consistency} lists the credit for each calculation stage of a transition-state step; a complete verified chain receives 1.0. The evaluator searches the same run's sandbox for the native files and links between calculations, recovering file associations where needed. It assigns credit from these results using the common physical criteria.

\begin{table}[htbp]
\centering
\caption{Graded self-consistency contribution for an accessible transition-state step. The endpoint contribution is 0.25 for one verified side and 0.50 for both sides.}
\label{tab:si-self-consistency}
\small
\begin{tabular}{p{0.65\textwidth}r}
\toprule
Evidence stage & Contribution \\
\midrule
Verified first-order saddle: exactly one imaginary frequency & 0.35 \\
IRC starting geometry consistent with the selected saddle & 0.15 \\
Native outer-frame--optimization--frequency lineage to both declared endpoint minima & 0.50 \\
\bottomrule
\end{tabular}
\end{table}

For graphs containing reaction edges, $C$ is the mean of the delivered edge scores. An edge represented by its endpoints alone requires coverage of both endpoint species. For a species-only graph, $C$ is the mean species-coverage score. Explicit inaccessible targets are exempt from requirements for stationary points or paths judged physically nonexistent under the task's permitted methods and receive the corresponding self-consistency credit.

The binary self-consistency check requires all requested species and steps to pass, including minimum characterization and the full transition-state/IRC chain. Counts of $C=1$ in the main text and score tables refer to full graded component scores. The binary check separately establishes whether the complete mechanism is validated. Fractional scores indicate which calculation stages have been supported.

\subsection{Electronic barriers and numerical accuracy}
For a step with a numerical reference, the evaluator checks the declared barrier against the transition-state and starting-species energies in the native outputs. If the barrier field is empty, it can calculate the difference from those energies when both are reported and supported by their native records. The conversion factor from hartree to kcal~mol$^{-1}$ is 627.509474.

For a supported numerical value $\widehat{\Delta E}_e^{\ddagger}$ and reference $\Delta E_{e,\mathrm{Gold}}^{\ddagger}$, the underlying Gaussian score is
\begin{equation}
q_e=\exp\left[-\frac{\left(\widehat{\Delta E}_e^{\ddagger}-
\Delta E_{e,\mathrm{Gold}}^{\ddagger}\right)^2}{2\sigma_e^2}\right].
\label{eq:si-gaussian}
\end{equation}
The released task rubrics use base widths of 3.0~kcal~mol$^{-1}$ for Easy and Impossible tasks and 2.0~kcal~mol$^{-1}$ for Normal and Hard tasks. For a reference barrier at or below 10.0~kcal~mol$^{-1}$, the reference-comparison width is the larger of the base width and 5.0~kcal~mol$^{-1}$. This small-barrier width applies to comparison with Gold; the declaration-to-log consistency check retains the task's base width.

All reference edges remain in the accuracy denominator with their recorded weights. A numerical reference edge requires a matching delivered step, a declared barrier, and a log-derived value to receive credit. Missing any of these gives zero for that edge; other supported steps keep their credit. The combined penalty for numerical deviations on supported edges is capped at 0.3. Penalties for missing steps, missing evidence, and incorrect negative conclusions apply in full. For total reference-edge weight $W$, supported numerically graded edges $D$, and missing or incorrectly delivered targets $U$,
\begin{equation}
A=1-\min\!\left(\frac{\sum_{e\in D}w_e(1-q_e)}{W},\,0.3\right)
-\frac{\sum_{e\in U}w_e}{W}.
\label{eq:si-accuracy-aggregation}
\end{equation}
Correct negative conclusions and accessible edges exempt from numerical references receive their prescribed full credit through $W$. For the latter, credit concerns the specified nonnumerical result.

The evaluator checks the allowed combinations of transition-state and reference energies in the delivered results. If one combination agrees with Gold and the reported barrier uses a different supported conformer or energy reference, the answer-definition rule replaces $q_e$ with 0.9. An unsupported barrier declaration with zero reference-consistent alternatives triggers the integrity rule: the self-consistency and accuracy contributions become zero, leaving $0.1G$. An omitted step loses its own contribution.

\subsection{Qualitative selectivity and property targets}
For HA1, HA6, HB1, and HF1, Accuracy scores the qualitative choice of favored channel. The delivered mechanism must identify one favored edge matching the reference channel or its product. Product identity is used when the agent's channel label differs from the reference label. Numerical barrier spreads are retained as calculations; the scored target is the qualitative channel preference. HL1 scores its independently referenced channel barriers.

For species-property tasks, Accuracy compares the declared properties with their references. Numerical comparison widths are 0.02~\AA\ for bond lengths, 50~cm$^{-1}$ for frequencies, and 2.0~kcal~mol$^{-1}$ for energy differences. State-ordering text is compared after case and whitespace normalization. When several properties are requested, their scores are averaged.

\subsection{Impossible tasks and final scientific claims}
Impossible tasks are scored on the agent's final conclusion about the requested target. Correctly identifying that target as inaccessible or beyond the allowed method's capability earns negative-target credit, including after exploratory calculations. Credit is lost for an incorrect claim that the target is accessible, a barrier or property presented as trustworthy for that negative target, or omission of the required negative conclusion. Endpoint species associated with a negative step receive the same treatment. A species-only negative target is evaluated directly on the species claim.

A result must answer the question about the requested electronic state, species, or pathway. The per-task tables and case analyses identify that target alongside its score.

\subsection{Original delivery, organized evaluation, and score reproduction}
Original delivery completeness records whether the run generated its mechanism graph during execution. Evaluation uses the graph organized by the common offline procedure in Section~\ref{sec:si-protocol}. The publication table records each run's final component scores and field sources. All configurations use the same final references and scoring rules.

The reproduction procedure reads the evaluation records to rebuild the per-run table, summaries, and figure data. It checks for 200 unique task/configuration records, 600 nonmissing component values, and agreement between the weighted components and the 200 final totals. The maximum absolute difference from the published four-decimal composites is $4.9711\times10^{-5}$. Rebuilding the publication tables uses the retained component records; independently reassessing the calculations also requires the run files and evaluator revision.
